\documentclass[12pt, a4paper]{article}

\usepackage[utf8]{inputenc}
\usepackage{graphicx}
\usepackage{hyperref}
\usepackage{float}
\usepackage{geometry}
\geometry{a4paper, margin=1in}
\usepackage{titlesec}
\usepackage{enumitem}
\usepackage{listings}
\usepackage{xcolor}

\definecolor{codegreen}{rgb}{0,0.6,0}
\definecolor{codegray}{rgb}{0.5,0.5,0.5}
\definecolor{codepurple}{rgb}{0.58,0,0.82}
\definecolor{backcolour}{rgb}{0.95,0.95,0.92}

\lstdefinestyle{mystyle}{
    backgroundcolor=\color{backcolour},   
    commentstyle=\color{codegreen},
    keywordstyle=\color{magenta},
    numberstyle=\tiny\color{codegray},
    stringstyle=\color{codepurple},
    basicstyle=\ttfamily\footnotesize,
    breakatwhitespace=false,         
    breaklines=true,                 
    captionpos=b,                    
    keepspaces=true,                 
    numbers=left,                    
    numbersep=5pt,                  
    showspaces=false,                
    showstringspaces=false,
    showtabs=false,                  
    tabsize=4
}
\lstset{style=mystyle}

\title{\textbf{DOKUMEN IMPLEMENTASI PROGRAM\\ Sistem Pengelolaan Nilai Akademik}}
\author{Dokumentasi Proyek}
\date{\today}

\begin{document}

\maketitle
\tableofcontents
\newpage

\section{Pendahuluan}
Dokumen ini merupakan laporan dari tahap \textbf{Implementasi Program} untuk Sistem Pengelolaan Nilai. Implementasi dibangun menggunakan bahasa pemrograman \textbf{PHP} dengan *framework* \textbf{Laravel}. Laporan ini menjabarkan pemetaan relasional, antarmuka fungsional, pencatatan *error*, dan pengaplikasian prinsip penulisan kode terstruktur serta berorientasi objek (OOP).

\section{Antarmuka dan Fungsionalitas Sistem}

\subsection{1. Halaman Login Berdasarkan Role}
Implementasi otentikasi login mendayagunakan kontrol kelas akses (*Authentication controller*) dan *Middleware* dengan pencabangan logika *enum* \texttt{role} pada tabel \texttt{users}. Setelah kredensial (*email* dan *password*) tervalidasi, sesi aplikasi mengatur arahan (*redirect*):
\begin{itemize}
    \item Administrator masuk ke dasbor \texttt{/admin/dashboard}.
    \item Guru diarahkan ke rute \texttt{/guru/dashboard} untuk operasi kelas harian.
    \item Siswa diarahkan ke \texttt{/siswa/dashboard} untuk meninjau rapor secara pasif.
\end{itemize}

\subsection{2. Form Input Data Siswa dan Nilai}
Halaman ini melayani fungsi CRUD melalui rute \texttt{/guru/input-nilai}. Dirancang menggunakan kerangka HTML berbalut kelas utilitas *TailwindCSS*. Form memuat *dropdown* \textit{Siswa} (\texttt{siswa\_id}), kolom *Mata Pelajaran*, dan tiga input desimal: Nilai Tugas, UTS, dan UAS. *Method action* diarahkan menuju ke \texttt{GuruController@storeNilai}.

\subsection{3. Proses Perhitungan Nilai Akhir}
Proses kalkulasi telah diabstraksi langsung ke dalam kelas Model \texttt{Nilai} untuk menghindari *Fat Controllers*. 
\begin{center}
\texttt{Nilai Akhir = (Tugas) + (UTS) + (UAS) / Total Rasio (diatur di logic method)}
\end{center}
Selain kalkulasi numerik desimal \texttt{5,2}, sistem juga memiliki prosedur \texttt{tentukanStatus()} yang mengatur atribut \texttt{status} menjadi \texttt{'Lulus'} atau \texttt{'Tidak Lulus'} secara otomatis sebelum tersimpan permanen di *database*.

\subsection{4. Laporan Hasil Nilai Siswa}
Dihasilkan melalui rute pencetakan (\texttt{/laporan/print}). Modul ini meng-kueri nilai berdasarkan identifikasi unik siswa (\texttt{nis}), dirender secara dinamis melalui mesin *Blade template* menjadi struktur laporan (Rapor), dan dapat dikonversi menjadi berkas PDF *offline*.

\section{Bukti Pengujian Database}
Pengujian ke \textit{database} MySQL dilakukan melalui eksekusi berkas migrasi (\textit{Migration Classes}). 
\begin{itemize}
    \item \textbf{Status Koneksi}: \textbf{SUKSES} (PDO tervalidasi stabil tanpa hambatan soket).
    \item \textbf{Bukti Tercipta Tabel:} Tersedia skema asli: \texttt{users}, \texttt{siswas}, \texttt{gurus}, dan \texttt{nilais}.
    \item Relasi (*Foreign Keys*), seperti penghapusan mengkaskade (\textit{cascade on delete}) untuk relasi \texttt{nilais.siswa\_id -> siswas.id}, berhasil terdefinisi utuh.
\end{itemize}

\section{Catatan Error (Debugging) \& Perbaikan}

\subsection{Error 1: \texttt{QueryException} pada Relasi Pengguna}
\textbf{Skenario:} Terjadi eror konstrain SQL (*foreign key constraint violation*) saat mencoba menghapus seorang pengguna (\texttt{user}) yang sedang terhubung dengan tabel data siswa. \\
\textbf{Perbaikan:} Memperbarui fail migrasi *schema* dengan menyisipkan aturan *constraint* yang lebih aman, misalnya \texttt{->nullOnDelete()} untuk rujukan \texttt{user\_id} pada tabel \texttt{siswas} dan \texttt{gurus}.
\begin{lstlisting}[language=PHP]
// Perbaikan pada Migrasi
$table->foreignId('user_id')->nullable()->constrained('users')->nullOnDelete();
\end{lstlisting}

\subsection{Error 2: Kalkulasi Status Tertinggal}
\textbf{Skenario:} Form guru hanya mengirim parameter Tugas, UTS, UAS. Nilai Akhir serta Status tidak ikut terisi ke *database* walau *store method* telah diaktifkan, memicu status \textit{null} atau \textit{0}. \\
\textbf{Perbaikan:} Memanfaatkan \textit{bootable saving model events} pada kelas OOP \texttt{Nilai} agar kedua nilai dikalkulasi mandiri persis sebelum penyimpanan DB dieksekusi.
\begin{lstlisting}[language=PHP]
// Pada model Nilai.php
protected static function booted(): void {
    static::saving(function (Nilai $nilai) {
        $nilai->nilai_akhir = $nilai->hitungNilaiAkhir();
        $nilai->status = $nilai->tentukanStatus();
    });
}
\end{lstlisting}

\section{Potongan Kode: Pemrograman Terstruktur \& OOP}

Sistem ini menunjukkan kapabilitas \textit{Object-Oriented Programming} yang rapi sesuai panduan internal proyek (*project architecture*):

\textbf{1. Implementasi Kelas Objek (\textit{Model}) \& Fungsi Terstruktur Internal}
Metode berikut membuktikan implementasi fungsi/prosedur yang membedah properti objek ke dalam *output* yang terstruktur.
\begin{lstlisting}[language=PHP]
<?php
namespace App\Models;
use Illuminate\Database\Eloquent\Model;

class Nilai extends Model
{
    // Fungsi prosedural / kalkulasi bisnis
    public function hitungNilaiAkhir(): float
    {
        // Contoh algoritma rasio rata-rata 33% 
        return round(($this->nilai_tugas + $this->nilai_uts + $this->nilai_uas) / 3, 2);
    }

    // Fungsi prosedural / klasifikasi kondisional
    public function tentukanStatus(): string
    {
        return $this->hitungNilaiAkhir() >= 75 ? 'Lulus' : 'Tidak Lulus';
    }
}
\end{lstlisting}

\textbf{2. Implementasi pada Kelas \textit{Controller}}
Implementasi struktur MVC Controller menangani validasi, otorisasi, dan perutean tanggapan (*response routing*).
\begin{lstlisting}[language=PHP]
<?php
namespace App\Http\Controllers;
use App\Models\Nilai;
use Illuminate\Http\Request;

class GuruController extends Controller
{
    // Method kelas untuk menyimpan data masukan guru
    public function storeNilai(Request $request)
    {
        $validatedData = $request->validate([
            'siswa_id' => 'required|exists:siswas,id',
            'mata_pelajaran' => 'required|string|max:255',
            'nilai_tugas' => 'required|numeric|min:0|max:100',
            'nilai_uts' => 'required|numeric|min:0|max:100',
            'nilai_uas' => 'required|numeric|min:0|max:100',
        ]);
        
        $nilai = new Nilai();
        $nilai->fill($validatedData);
        $nilai->guru_id = auth()->user()->guru->id;
        $nilai->save(); // Method saving otomatis menghitung nilai akhir & status (Event Hook)
        
        return redirect()->route('guru.rekap')->with('success', 'Nilai tersimpan.');
    }
}
\end{lstlisting}

\section{Pemanfaatan Library dan Komponen}
Aplikasi ini ditunjang dengan komponen standar:
\begin{enumerate}
    \item \textbf{Laravel 11.x:} Inti *framework* PHP moderen yang mendukung ekosistem MVC secara penuh dan aman (proteksi CSRF terintegrasi).
    \item \textbf{Tailwind CSS:} Paket kompiler kelas bergaya atomik (\textit{utility-first CSS}) yang memodernisasi desain tata letak UI tanpa harus menulis kustomisasi kode CSS secara masif.
    \item \textbf{Eloquent ORM:} Lapisan pemetaan relasional bawaan Laravel yang memungkinkan interaksi *database* secara *fluent* dan objektif tanpa instruksi sintaks SQL manual.
\end{enumerate}

\section{Penjelasan Coding Guidelines dan Best Practices}
Sistem Pengelolaan Nilai diprogram berdasarkan standar emas industri kode program:
\begin{itemize}
    \item \textbf{Fat Models, Skinny Controllers:} Kontroler seperti \texttt{GuruController} dipertahankan ringkas. Segala rumusan bisnis (\texttt{hitungNilaiAkhir()}) dipindahkan langsung ke berkas representasi Model \texttt{Nilai}.
    \item \textbf{Penggunaan Model Event Observers:} Memanfaatkan pemicu kejadian sistematis (\texttt{static::saving}) ketimbang harus memanggil secara manual formula perhitungan tiap kali kontroler melakukan simpan (*DRY principle - Don't Repeat Yourself*).
    \item \textbf{Penamaan Metrik Kaku (Convention over Configuration):} Format penamaan singular pada Model (\texttt{Siswa}, \texttt{Guru}), dan format *plural* pada *database migrations* (\texttt{siswas}, \texttt{gurus}) dilakukan secara konsisten agar Eloquent dapat menterjemahkannya secara sihir (*magically resolved*).
\end{itemize}

\end{document}
